\scanpage{176}

U N I V E R S I T E

DES SCIENCES ET TECHNIQUES

DU LANGUEDOC

\verb|^^^^^^^^^^^^|

\hfill Montpellier, le 14 octobre 1980

Service de la Scolarité

AM/gs

Diffusion : Liste A

N° 5256

\begin{center}
\uline{AVIS DE SOUTENANCE D'UNE THESE DE 3e CYCLE}

``\uline{PHYSIQUE THEORIQUE option ASTROPHYSIQUE}''
\end{center}

\uline{NOM et prénom du candidat} : GRAL Joëlle

\uline{Titre de la thèse} : ``ETUDE DES ETOILES VARIABLES A LONGUE PERIODE PROPRIETES DES ONDES DE CHOC ET ENVELOPPES CIRCUMSTELLAIRES''

\uline{Date, heure et lieu de soutenance prévus} : mardi 21 octobre 1980 à 16 h.,30

Bâtiment Mathématiques - Université des Sciences et Techniques du Languedoc.

\uline{Nom du Directeur de thèse} : M. H. ANDRILLAT

\uline{Désignation du Laboratoire où la thèse est préparée} : Laboratoire d'Astronomie

\hfill U.S.T.L.

\uline{Analyse de l'ouvrage} : Le spectre des étoiles variables à longue période, et plus particulièrement les raies d'émission peuvent être expliquées par le passage d'une onde de choc dans l'atmosphère de l'étoile. Après quelques géné-\newline
ralités sur les ondes de choc, je me suis penchée sur la vitesse de l'onde ; elle est assez grande pour qu'une enveloppe de matière soit éjectée à chaque dilata-\newline
tion de l'étoile. J'ai donc abordé le problème de perte de masse et l'étude des enveloppes circumstellaires ce qui conduit à donner plusieurs modèles statiques et dynamiques sur la structure des étoiles variables.

\uline{Durée approximative des recherches ayant abouti à la thèse} : 2 ans

\uline{Composition du jury} :
\begin{tabular}{l}
M. ANDRILLAT H.\\
M. SABATIER P.\\
M. DELOUPY Cl.\\
Mme MENNESSIER M.O.
\end{tabular}

\begin{flushright}
P/ le Secrétaire Général,\\
l'Attachée d'administration\\
Universitaire,

\unsure{Gentil}

\uline{A. GENTIL}
\end{flushright}
\margin{176}
